\documentclass[aspectratio=169,10pt,spanish]{beamer}

\input{../../../_preambulo_beamer.tex}
\input{../../../_comandos_beamer.tex}

\selectlanguage{spanish}

\title{MA1001B - Modelación Estadística para la Toma de Decisiones}
\subtitle{Lección 32: Distribución F y comparación de dos varianzas}
\author{}
\institute{Tecnol\'ogico de Monterrey}
\date{}

\begin{document}

\begin{frame}[plain]
  \vspace{1.2cm}
  {\Large\bfseries MA1001B - Modelación Estadística para la Toma de Decisiones\par}
  \vspace{0.7cm}
  {\large Lección 32: Distribución F y comparación de dos varianzas\par}
  \vspace{0.7cm}
  {\color{TecRojo}\rule{\textwidth}{1.2pt}}
  \vspace{0.8cm}

  Facultad MA1001B\\[0.3cm]
  Periodo\\[0.3cm]
  Tecnol\'ogico de Monterrey
\end{frame}

\begin{frame}{La pregunta de las dos varianzas}
  \begin{alertblock}{Pregunta guía}
    ¿Cómo se comparan dos varianzas muestrales independientes mediante una
    razón $F$ y una probabilidad de cola superior?
  \end{alertblock}

  \vspace{0.6em}
  Al final de esta lección, usted puede:
  \begin{itemize}
    \item calcular $s_{1}^{2}$ y $s_{2}^{2}$ a partir de dos muestras;
    \item formar $F=s_{1}^{2}/s_{2}^{2}$ con un par de grados de libertad;
    \item leer $p$ de \texttt{scipy.stats.f.sf};
    \item explicar por qué la prueba $F$ de dos varianzas es sensible a la no
      normalidad.
  \end{itemize}

  \vfill
  \scriptsize Todos los tiempos de ciclo usados hoy son sintéticos. No se usan datos reales de empresa.
\end{frame}

\begin{frame}{Dos estaciones de empaque independientes}
  \small
  \begin{center}
    \begin{tabular}{lrr}
      \toprule
      \textbf{Cantidad} & \textbf{Estación A} & \textbf{Estación B} \\
      \midrule
      Tamaño de muestra & $n_{1}=25$ & $n_{2}=25$ \\
      Grados de libertad & 24 & 24 \\
      Varianza muestral & $s_{1}^{2}=44.205731$ & $s_{2}^{2}=12.024693$ \\
      \bottomrule
    \end{tabular}
  \end{center}

  \vspace{0.5em}
  \begin{exampleblock}{Hipótesis}
    $H_{0}:\sigma_{1}^{2}=\sigma_{2}^{2}$ frente a $H_{1}:\sigma_{1}^{2}>\sigma_{2}^{2}$.
    Las estaciones son independientes. Las medias no son el objetivo de esta prueba.
  \end{exampleblock}
\end{frame}

\begin{frame}[fragile]{Paso 1: Dos varianzas muestrales}
  \begin{columns}[T]
    \begin{column}{0.42\textwidth}
      \small
      \begin{lstlisting}
s1_sq = np.var(a, ddof=1)
s2_sq = np.var(b, ddof=1)
      \end{lstlisting}

      \begin{block}{Salida verificada}
        $s_{1}^{2}=44.205731$\\
        $s_{2}^{2}=12.024693$\\
        $\mathrm{gl}=(24,24)$
      \end{block}
    \end{column}
    \begin{column}{0.58\textwidth}
      \centering
      \includegraphics[width=\textwidth]{../figuras/f_dos_var_01_varianzas_muestrales.png}
      \vspace{0.2em}
      \scriptsize Diagramas de caja de las estaciones del Paso 1
    \end{column}
  \end{columns}
\end{frame}

\begin{frame}{La razón $F$ y la cola superior}
  \small
  \[
    F=\frac{s_{1}^{2}}{s_{2}^{2}}=\frac{44.205731}{12.024693}=3.676246.
  \]
  Valor crítico y valor $p$:
  \[
    F_{0.05}(24,24)=1.983760,
  \]
  \[
    p=\texttt{f.sf}(3.676246,24,24)=0.001125.
  \]
  Como $p<\alpha=0.05$, se rechaza $H_{0}$: la estación A es más variable
  en esta instantánea sintética.
\end{frame}

\begin{frame}[fragile]{Paso 2: $F$ y \texttt{f.sf}}
  \begin{columns}[T]
    \begin{column}{0.40\textwidth}
      \small
      \begin{lstlisting}
F = s1_sq / s2_sq
p = stats.f.sf(F, 24, 24)
      \end{lstlisting}

      \begin{block}{Salida verificada}
        $F=3.676246$\\
        $p=0.001125$\\
        Rechazar $H_{0}$
      \end{block}
    \end{column}
    \begin{column}{0.60\textwidth}
      \centering
      \includegraphics[width=\textwidth]{../figuras/f_dos_var_02_prueba_f.png}
      \vspace{0.2em}
      \scriptsize Cola superior $F$ del Paso 2
    \end{column}
  \end{columns}
\end{frame}

\begin{frame}{Paso 3: $F$ es sensible a la no normalidad}
  \begin{columns}[T]
    \begin{column}{0.42\textwidth}
      \small
      Varianzas iguales, 5000 pruebas, $\alpha=0.05$.

      \vspace{0.4em}
      Muestras normales: tasa de tipo I $0.052200$.

      \vspace{0.4em}
      Muestras $t(3)$: tasa de tipo I $0.202400$.

      \vspace{0.5em}
      \textbf{Límite:} las colas pesadas inflan los rechazos falsos aun cuando
      las dos varianzas son iguales. La prueba $F$ de dos varianzas no es
      robusta.
    \end{column}
    \begin{column}{0.58\textwidth}
      \centering
      \includegraphics[width=\textwidth]{../figuras/f_dos_var_03_sensibilidad_no_normal.png}
      \vspace{0.2em}
      \scriptsize Tasas de tipo I del Paso 3
    \end{column}
  \end{columns}
\end{frame}

\begin{frame}{Verificación de conceptos}
  \begin{enumerate}
    \item Calcule $F$ a partir de $s_{1}^{2}=44.205731$ y $s_{2}^{2}=12.024693$.
    \item ¿Cuáles son los dos grados de libertad para $n_{1}=n_{2}=25$?
    \item ¿$p=0.001125$ compara medias o varianzas?
    \item ¿Por qué una tasa de tipo I de $0.202400$ es un problema con
      $\alpha=0.05$?
  \end{enumerate}

  \vspace{0.6em}
  \begin{block}{Conexión con la Lección 33}
    La sesión puede continuar directamente con $H_{0}$, $H_{1}$, errores de
    tipo I y tipo II como marco general de pruebas.
  \end{block}

  \vfill
  \scriptsize Fuente: OpenStax, \textit{Introductory Business Statistics 2e},
  Capítulo 12, Sección 12.1. CC BY-NC-SA.
\end{frame}

\appendix



\end{document}
